\scanpage{205}
% archive sheet number: 204
% top, continuation from scan 204; the first deleted word stands above the second
classe de faisceaux inversibles \deleted{\unsure{de}} \deleted{\unsure{locaux}} sur $S$ est dans le Groupe de Brauer $H^2(S,\unsure{\mathbb G_m|G_m})$ [\unsure{se rappeler|se rappelle} que
\[
\begin{aligned}
H^*(S,\mathbb G_m)&\underset{p}{\Longleftarrow} H^p(\unsure{B|S},\underline H^q_{\unsure{\mathrm{local}|\mathrm{loc}}}(\mathbb G_m))\quad\text{donc}\\
H^n(S,\mathbb G_m)&\unsure{\longrightarrow|-} H^0(S,\underline H^n_{\unsure{\mathrm{loc}|\mathrm{local}}}(\mathbb G_m))\quad\cdots
\end{aligned}
\]
]. On obtient ainsi l'\uline{accouplement de Tate}
\[
\boxed{H^1(S,A)\times H^0(S,A')\longrightarrow H^2(S,\mathbb G_m)}
\]

% middle
\unsure{Ceci} d'ailleurs se \unsure{précise}. (\unsure{Supposons p.ex.} que $S'\longrightarrow S$ est fid plat) \unsure{on a} des accouplements
\[
H^1(S'/S,A)\times H^0(S,A')\longrightarrow H^2(S'/S,\mathbb G_m).
\]

% bottom
Il semble, quand $S'$ est \uline{fini} sur $S$, de sorte que les $\widehat H^i(S'/S,A)$ sont définis, qu'on puisse définir \unsure{même} des accouplements
\[
\widehat H^i(S'/S,A)\times\widehat H^{1-i}(S'/S,A')\longrightarrow H^2(S'/S,\mathbb G_m)
\]
qu'il faudrait comprendre \ldots\ldots
